\section{Experimental Evaluation}
\label{sec:evaluation}

\subsection{Comparison Setup}
\label{sec:experimental-setup}

We compare GA-LK with ZAC\cite{lin2025zac} and Stade et al.'s routing-aware placement\cite{stade2025routingaware}, using their published implementations.\footnote{ZAC uses ASAP scheduling with dynamic placement and reuse. Routing-aware placement uses MQT QMAP 3.2.0/Core 3.1.0, with look-ahead factor 0.2, deepening value 0.2, reuse level 5, and deepening factor 0.6.} Inputs comprise 18 ZAC benchmarks (ZAC18)\cite{lin2025zac} and 154 files from the MQT QMAP example set\cite{wille2023qmap}, representing 151 distinct canonicalized circuits. Equivalent inputs are merged for aggregation. All compilers receive the same gate sequences. We evaluate the baselines' original schedules under their recorded acceptance criteria, without additional repair or diagnostic filtering. GA-LK checks non-target intersections for conflicts with stationary atoms during candidate evaluation.

The shared architecture has one AOD, $100\times100$ storage traps, and two paired $7\times20$ entangling arrays providing 140 gate sites. In~\eqref{eq:fidelity}, $(f_1,\allowbreak f_2,\allowbreak f_{\rm tran},\allowbreak f_{\rm exc})=(0.9997,\allowbreak 0.995,\allowbreak 0.999,\allowbreak 0.9975)$ and $T_2=1.5$\,s. Single-qubit gates, CZ gates, and trap transfers take 52, 0.36, and 15\,$\mu$s, respectively. Quality comparisons include circuits compiled by every participating configuration with each atom's idle time below $T_2$. The main experiment requires three valid GA-LK seeds and takes their per-circuit median. Across circuits, fidelity uses geometric means; batches and rearrangement latency use arithmetic means.

The main comparison uses each compiler's initialization. GA-LK evaluates at most 32 mappings with population 8 over the first layer and the next two. Subsequent layers within this prefix permit 32 encoding evaluations with inner look-ahead disabled. Dynamic optimization uses $H=8$, $\alpha=0.5$, $\rho=0.7$, $\epsilon=0.05$, and at most 576 encoding evaluations, reduced with circuit size. Look-ahead, search, and budget comparisons fix initial mappings; initialization comparisons fix subsequent dynamic optimization. Each study uses the sample defined by its compared configurations.

\subsection{Overall Gains and Physical Sources}

GA-LK improves geometric-mean fidelity and reduces mean rearrangement batches on both circuit sets against both baselines, as shown in Table~\ref{tab:main-results}. Relative to ZAC, fidelity increases by \GAMainZACVsZACFidelityGain\% on ZAC18 and \GAMainQMAPVsZACFidelityGain\% on QMAP; QMAP batches decrease by \GAMainQMAPVsZACBatchReduction\%.

\begin{table}[!tb]
\caption{\PaperRevision{Overall results for three compilers on two circuit sets}}
\label{tab:main-results}
\centering
\PaperTableSetup
\fontsize{8.5pt}{9.7pt}\selectfont
\setlength{\tabcolsep}{1.2pt}
\begin{tabular*}{\columnwidth}{@{\extracolsep{\fill}}llrrrr@{}}
\toprule
Set & Method & \PaperRevision{Fidelity} & \PaperRevision{Batches} & \shortstack{\PaperRevision{Rearr. time}\\\PaperRevision{(ms)}} & \shortstack{\PaperRevision{Compiled}\\\PaperRevision{/inputs}}\\
\midrule
\multirow{3}{*}{\shortstack{ZAC18\\\PaperRevision{$N=\GAMainZACStrictN$}}}
& ZAC & \PaperRevision{\pgfmathprintnumber[fixed,precision=4,zerofill]{\GAMainZACMOneF}} & \PaperRevision{\GAMainZACMOneB} & \PaperRevision{\GAMainZACMOneT} & \PaperRevision{\GAMainZACMOneV}\\
& Stade et al. & \PaperRevision{\pgfmathprintnumber[fixed,precision=4,zerofill]{\GAMainZACMTwoF}} & \PaperRevision{\GAMainZACMTwoB} & \PaperRevision{\textbf{\GAMainZACMTwoT}} & \PaperRevision{\GAMainZACMTwoV}\\
& GA-LK & \PaperRevision{\textbf{\boldmath\pgfmathprintnumber[fixed,precision=4,zerofill]{\GAMainZACMFourF}}} & \PaperRevision{\textbf{\GAMainZACMFourB}} & \PaperRevision{\GAMainZACMFourT} & \PaperRevision{\GAMainZACMFourV}\\
\midrule
\multirow{3}{*}{\shortstack{\PaperRevision{QMAP}\\\PaperRevision{$N=\GAMainQMAPStrictN$}}}
& ZAC & \PaperRevision{\pgfmathprintnumber[fixed,precision=6,zerofill]{\GAMainQMAPMOneF}} & \PaperRevision{\GAMainQMAPMOneB} & \PaperRevision{\GAMainQMAPMOneT} & \PaperRevision{\GAMainQMAPMOneV}\\
& Stade et al. & \PaperRevision{\pgfmathprintnumber[fixed,precision=6,zerofill]{\GAMainQMAPMTwoF}} & \PaperRevision{\GAMainQMAPMTwoB} & \PaperRevision{\GAMainQMAPMTwoT} & \PaperRevision{\GAMainQMAPMTwoV}\\
& GA-LK & \PaperRevision{\textbf{\boldmath\pgfmathprintnumber[fixed,precision=6,zerofill]{\GAMainQMAPMFourF}}} & \PaperRevision{\textbf{\GAMainQMAPMFourB}} & \PaperRevision{\textbf{\GAMainQMAPMFourT}} & \PaperRevision{\GAMainQMAPMFourV}\\
\bottomrule
\end{tabular*}
\PaperTableNote[\columnwidth]{\fontsize{8pt}{9.1pt}\selectfont\PaperRevision{$N$ counts distinct circuits included in the quality comparison. Fidelity uses the geometric mean; rearrangement batches and time use arithmetic means. The \GAMainQMAPStrictFileN{} QMAP files represent \GAMainQMAPStrictN{} distinct circuits. The last column counts completed compilations over all input files, including results outside the fidelity model's domain. Bold indicates the best value within each circuit set.}}
\end{table}


The gains reflect different circuit-level patterns. Against ZAC, GA-LK improves \GAMainZACVsZACWins{} ZAC18 circuits and lowers fidelity on \GAMainZACVsZACLosses{}; the QMAP counts are \GAMainQMAPVsZACWins{} and \GAMainQMAPVsZACLosses{}, respectively. Against the routing-aware baseline, QMAP yields \GAMainQMAPVsICCADWins{} wins and \GAMainQMAPVsICCADLosses{} losses. Thus, the QMAP aggregate reflects larger gains on a subset of circuits, rather than more frequent wins. The geometric mean captures this balance through the average log fidelity ratio.

Across both circuit sets and baselines, transfer and decoherence savings outweigh added idle-excitation losses. For QFT-18, transfers fall from \GAMainMechanismQFTBaseTransfers{} with ZAC to \GAMainMechanismQFTGATransfers{}; seed-median contributions to $\Delta\ln F$ from transfer, excitation, and decoherence are \mbox{$(+\GAMainMechanismQFTTransferGainRounded,\,\GAMainMechanismQFTExcitationGainRounded,\,+\GAMainMechanismQFTCoherenceGainRounded)$}. Transfer savings dominate this case.

\subsection{Effects of Look-Ahead and Search}
\label{sec:ablation}

Within the same joint decision space, looking beyond the current layer improves plan selection. With initial mappings, joint variables, search budgets, and candidate storage domains fixed, changing $H=0$ to $H=8$ raises geometric-mean fidelity by \PaperPercentGain{\AblationHRatio}{1} across \AblationHN{} circuits. The $H=8$ configuration accounts for subsequent gates, terminal returns, and residency informed by near-term reuse; $H=0$ selects by current-layer loss alone.

A separate comparison fixes the cross-layer objective to assess how genetic search selects among joint candidates. Above the enumeration threshold, genetic and coordinate-wise greedy search use the same $H=8$, constraints, objective, evaluation budget, and local-refinement rules. Genetic search improves geometric-mean fidelity by \PaperPercentGain{\AblationGARatio}{1} across \AblationGAN{} circuits. These comparisons separately assess look-ahead configuration and search within the same joint space. They use different samples, so their increments cannot be added.

\label{sec:horizon-evaluation}
The value of look-ahead also depends on its range. In a separate horizon study, most of the observed average improvement is already obtained with one-layer look-ahead. Fig.~\ref{fig:experimental-summary} shows the complete trajectories of \HorizonExtCompleteN{} circuits: moving from $H=0$ to 1 raises geometric-mean fidelity ratios and lowers geometric-mean batch ratios. Larger bounds produce smaller average changes. Four circuits have identical actual windows at $H=4$ and 8, reflecting the remaining-layer and size limits on prediction.

\begin{figure}[!t]
  \centering
  \resizebox{\columnwidth}{!}{% Two side-by-side plots at final single-column width, following the supplied template.
% All ten circuit medians and both geometric means retain the frozen H8 normalization.
\begingroup
\PaperShowRevisionsfalse
% Shared vector-figure vocabulary for the Chinese IEEE manuscript.
% Every figure in this directory inputs this file, so the manuscript only
% needs the standard TikZ package and the libraries already loaded by
% paper_zh.tex.  Keep the guard: figures may be input more than once while
% editing or when a stand-alone smoke-test document is built.
\ifdefined\GALKFigureStyleLoaded\else
\def\GALKFigureStyleLoaded{1}

% Indigo / blue / ochre palette adapted from Paul Tol qualitative schemes.
% Source: https://sronpersonalpages.nl/~pault/ (accessed 2026-09-14).
% Dark strokes identify roles; pale fills preserve text contrast.
% Atoms are filled, traps hollow, predictions dashed, invalid paths crossed.
\definecolor{galkInk}{HTML}{263238}
\definecolor{galkMuted}{HTML}{727C86}
\definecolor{galkGrid}{HTML}{E0E5EA}
\definecolor{galkFrame}{HTML}{A4AFBA}
\definecolor{galkPaper}{HTML}{FFFFFF}
% Blue and ochre identify storage and entangling zones; indigo highlights GA-LK.
% Forecasts use magenta dashes; comparison series retain grey open markers.
\definecolor{galkStorage}{HTML}{4477AA}
\definecolor{galkStorageFill}{HTML}{EDF4FB}
\definecolor{galkEntangle}{HTML}{997700}
\definecolor{galkEntangleFill}{HTML}{FFF6D9}
\definecolor{galkResident}{HTML}{332288}
\definecolor{galkResidentFill}{HTML}{F0EFF8}
\definecolor{galkGood}{HTML}{263238}
\definecolor{galkBad}{HTML}{CC3311}
\definecolor{galkLook}{HTML}{AA3377}
\definecolor{galkLookFill}{HTML}{F7EDF3}
\definecolor{galkData}{HTML}{332288}
\definecolor{galkTail}{HTML}{929AA3}
\definecolor{galkZAC}{HTML}{727C86}
\definecolor{galkICCAD}{HTML}{929AA3}
\definecolor{galkGroupOne}{HTML}{332288}
\definecolor{galkGroupTwo}{HTML}{727C86}
\definecolor{galkGroupThree}{HTML}{929AA3}

\tikzset{
  % Shared engineering-figure contract.  Physical snapshots must place an
  % paired-trap entanglement band above a dense storage field; only their scale may
  % change.  Labels are aligned at the upper-left of both regions.
  galk/base/.style={
    font=\rmfamily\fontsize{9.0pt}{10.0pt}\selectfont,
    text=galkInk,
    line cap=butt,
    line join=miter,
    >=Stealth
  },
  galk/panel title/.style={
    font=\rmfamily\bfseries\fontsize{9.0pt}{10.0pt}\selectfont,
    text=galkInk,
    anchor=west
  },
  galk/label/.style={
    font=\rmfamily\fontsize{9.0pt}{10.0pt}\selectfont,
    text=galkInk
  },
  galk/small label/.style={
    font=\rmfamily\fontsize{8.7pt}{9.7pt}\selectfont,
    text=galkInk
  },
  galk/zone label/.style={
    font=\rmfamily\bfseries\fontsize{8.7pt}{9.7pt}\selectfont,
    text=galkInk,
    anchor=north west,
    inner sep=.5pt
  },
  galk/entanglement zone/.style={
    draw=galkEntangle,
    fill=galkEntangleFill,
    line width=.65pt
  },
  galk/storage zone/.style={
    draw=galkStorage,
    fill=galkStorageFill,
    line width=.65pt
  },
  galk/site/.style={
    circle,
    draw=galkInk!82,
    fill=white,
    line width=.36pt,
    minimum size=1.35mm,
    inner sep=0pt
  },
  galk/occupied atom/.style={
    circle,
    draw=galkInk,
    fill=galkInk,
    line width=.36pt,
    minimum size=1.55mm,
    inner sep=0pt
  },
  galk/target site/.style={
    circle,
    draw=galkInk!60,
    fill=galkInk!10,
    line width=.34pt,
    minimum size=1.55mm,
    inner sep=0pt
  },
  galk/q label/.style={
    font=\rmfamily\itshape\fontsize{8.5pt}{9.5pt}\selectfont,
    text=galkInk,
    inner sep=.25pt
  },
  galk/rydberg pair/.style={
    draw=galkInk!72,
    densely dashed,
    line width=.36pt
  },
  galk/accepted path/.style={
    -{Stealth[length=.95mm,width=.62mm]},
    draw=galkInk,
    line width=.70pt
  },
  galk/alternative path/.style={
    -{Stealth[length=.95mm,width=.62mm]},
    draw=galkMuted,
    densely dashed,
    line width=.65pt
  },
  galk/future path/.style={
    -{Stealth[length=.95mm,width=.62mm]},
    draw=galkLook,
    densely dashed,
    line width=.70pt
  },
  galk/invalid path/.style={
    -{Stealth[length=.95mm,width=.62mm]},
    draw=galkBad,
    densely dashed,
    line width=.70pt
  }
}
\fi

\pgfplotstableread[col sep=space]{figures/data/horizon_trends.dat}\HorizonTrendData
\begin{tikzpicture}[galk/base,
  horizon label/.style={font=\rmfamily\fontsize{10pt}{11pt}\selectfont}]
  \path[use as bounding box] (0,0) rectangle (8.80,4.49);
  \node[horizon label,font=\rmfamily\bfseries\fontsize{10pt}{11pt}\selectfont]
    at (2.70,4.24) {Fidelity};
  \node[horizon label,font=\rmfamily\bfseries\fontsize{10pt}{11pt}\selectfont]
    at (6.73,4.24) {Batches};
  \pgfplotsset{horizon trend axis/.style={
    width=3.16cm,height=2.45cm,scale only axis,
    axis lines*=left,axis line style={draw=galkInk,line width=.65pt},
    tick style={draw=galkInk,line width=.45pt},tick align=outside,
    tick label style={font=\rmfamily\fontsize{10pt}{11pt}\selectfont},
    label style={font=\rmfamily\fontsize{10pt}{11pt}\selectfont},
    xmin=-.15,xmax=8.2,xtick={0,1,2,4,8},
    ymajorgrids=true,grid style={draw=galkMuted!28,line width=.35pt},
    scaled ticks=false,clip=true,
  }}
  \begin{axis}[horizon trend axis,
    at={(1.12cm,1.43cm)},anchor=south west,
    ylabel={Ratio to $H=8$},
    ylabel style={at={(axis description cs:-.27,.5)},anchor=south,inner sep=0pt},
    ymin=.18,ymax=1.055,ytick={.2,.4,.6,.8,1.0}]
    \addplot[draw=galkMuted!85,densely dashed,line width=.55pt,forget plot] coordinates {(0,1) (8,1)};
    \pgfplotsinvokeforeach{0,...,9}{
      \addplot[draw=galkMuted!70,line width=.65pt,mark=*,mark size=.95pt,
        mark options={draw=galkMuted!85,fill=white},forget plot]
        table[x=horizon,y=f#1]{\HorizonTrendData};
    }
    \addplot[draw=galkData,line width=1.3pt,mark=diamond*,mark size=2pt,
      mark options={draw=white,line width=.4pt,fill=galkData}]
      table[x=horizon,y=f_mean]{\HorizonTrendData};
  \end{axis}
  \begin{axis}[horizon trend axis,
    at={(5.15cm,1.43cm)},anchor=south west,
    ymin=.95,ymax=1.73,ytick={1.0,1.2,1.4,1.6}]
    \addplot[draw=galkMuted!85,densely dashed,line width=.55pt,forget plot] coordinates {(0,1) (8,1)};
    \pgfplotsinvokeforeach{0,...,9}{
      \addplot[draw=galkMuted!70,line width=.65pt,mark=*,mark size=.95pt,
        mark options={draw=galkMuted!85,fill=white},forget plot]
        table[x=horizon,y=b#1]{\HorizonTrendData};
    }
    \addplot[draw=galkData,line width=1.3pt,mark=diamond*,mark size=2pt,
      mark options={draw=white,line width=.4pt,fill=galkData}]
      table[x=horizon,y=b_mean]{\HorizonTrendData};
  \end{axis}
  \node[horizon label] at (4.65,.64) {Look-ahead upper bound $H$};
  \draw[draw=galkMuted!70,line width=.65pt] (.32,.20)--(.86,.20);
  \node[circle,draw=galkMuted!85,fill=white,inner sep=1pt] at (.59,.20) {};
  \node[horizon label,anchor=west] at (.96,.20) {Circuits ($N=10$)};
  \draw[draw=galkData,line width=1.3pt] (4.70,.20)--(5.24,.20);
  \node[diamond,draw=white,fill=galkData,inner sep=1.3pt] at (4.97,.20) {};
  \node[horizon label,anchor=west] at (5.34,.20) {Geometric mean};
\end{tikzpicture}
\endgroup
}
  \caption{Effect of the look-ahead bound. Each circuit uses the median of \HorizonExtSeedN{} seeds, normalized to its own $H=8$ result. Left and right plots show fidelity and batch ratios. Thin lines retain every circuit; thick lines show the geometric mean of \HorizonExtCompleteN{} circuits. Actual depth also depends on remaining layers and size budgets.}
  \label{fig:experimental-summary}
\end{figure}


\label{sec:init-evaluation}
The same principle applies before execution begins, when the initial layout sets the transport origins for early gates. Across \PhysicalInitSearchCompleteN{} circuits, evaluating the first layer and the next two improves full-circuit geometric-mean fidelity by \PhysicalInitGaVsHZeroGainPercent\% over first-layer-only evaluation. The comparison yields \PhysicalInitGaVsHZeroWins{} wins and \PhysicalInitGaVsHZeroLosses{} losses, using per-circuit medians over three seeds. With site domains, initial populations, prefix ranges, and 32-mapping budgets fixed, GA improves fidelity by \PhysicalInitGaVsRandomGainPercent\% over random search (\PhysicalInitGaVsRandomWins{} wins, \PhysicalInitGaVsRandomTies{} tie, \PhysicalInitGaVsRandomLosses{} losses). Both use dynamic optimization at $H=8$. Early-layer physical evaluation therefore informs initialization as well as dynamic placement, while genetic search improves the mapping choice under a matched budget.

\subsection{Computational Budget}
\label{sec:runtime-evaluation}

These quality gains require additional candidate evaluation, whose budget provides a direct way to control compilation cost. On an arm64 macOS host with 14 logical processors and 24\,GiB RAM, 127 circuits in the main quality study were compiled with up to 14 concurrent workers. Taking three-seed medians per circuit and then the median across circuits gives 19.12\,s for complete compilation, including GA initialization; initialization alone takes a median of 4.92\,s.

A separate serial study measures the cost after fixing initial layouts. Across \GAMainFixedLayoutCircuitN{} circuits, non-initialization compilation takes a median of \GAMainPostInitialMedianSeconds\,s on the same host. It uses identical initial mappings and seed 0, with three repetitions after \mbox{warm-up}. The statistic subtracts the entire initialization stage from each run, then takes medians within and across circuits.

With initial layouts fixed, smaller candidate sets lower compilation cost while retaining similar execution quality. Reducing encoding evaluations from 576 to 192, or candidate sites and assignments from $(6,4)$ to $(4,2)$, gives fidelity ratios close to 1 across \SensitivityFidelityN{} circuits. Across \GAMainBudgetTimingN{} timed circuits, compilation time excluding initialization falls by \PaperPercentReduction{\GAMainBudgetPostInitialTimeRatio}{1} and \PaperPercentReduction{\GAMainReturnPostInitialTimeRatio}{1}, respectively, based on geometric means of paired time ratios. Each configuration runs once with seed 0 and up to four concurrent runs; these reductions describe that schedule.
